%=====================================================================
%  Swayam Singh — Résumé
%  Custom single-column template (Overleaf-ready, compile with pdflatex)
%  Style: centered name, FontAwesome contact row, ruled section headers,
%         bold title / right-aligned date entries, tight bullet lists.
%=====================================================================
\documentclass[a4paper,9pt]{extarticle}

\usepackage[T1]{fontenc}
\usepackage[default]{lato}            % clean humanist sans (matches the original)
\usepackage[hidelinks]{hyperref}
\usepackage{fontawesome5}             % contact + inline icons
\usepackage{titlesec}
\usepackage{enumitem}
\usepackage[usenames,dvipsnames]{xcolor}
\usepackage{microtype}
\usepackage[
  a4paper,
  left=1.25cm, right=1.25cm, top=0.85cm, bottom=0.8cm
]{geometry}

% ---- global tweaks --------------------------------------------------
\setlength{\parindent}{0pt}
\setlength{\emergencystretch}{2em}     % avoid overfull lines into the margin
\linespread{0.97}
\pagestyle{empty}
\definecolor{accent}{HTML}{1A1A1A}     % near-black; change for a colored theme
\hypersetup{urlcolor=accent}

% ---- section headers: ALL CAPS, bold, full-width rule ---------------
\titleformat{\section}
  {\large\bfseries\color{accent}}
  {}{0em}{\MakeUppercase}
  [\vspace{-5pt}{\color{accent}\rule{\linewidth}{0.8pt}}\vspace{-6pt}]
\titlespacing*{\section}{0pt}{5pt}{3pt}

% ---- tight bullet list ----------------------------------------------
\newlist{tight}{itemize}{2}
\setlist[tight,1]{leftmargin=1.35em, itemsep=0.8pt, topsep=1.2pt, parsep=0pt,
                  label=\small$\bullet$}
\setlist[tight,2]{leftmargin=1.25em, itemsep=1pt,   topsep=1pt, parsep=0pt,
                  label=\small$\circ$}

% ---- helper: title (left, bold) + meta (right) ----------------------
\newcommand{\headerrow}[2]{\textbf{#1}\hfill #2\par}
\newcommand{\subrow}[2]{\textbf{#1}\hfill #2\par}

% ---- helper: an experience/education block --------------------------
% \entry{Title}{Right meta (dates)}{Organisation}{Right meta (location)}
\newcommand{\entry}[4]{%
  \vspace{1pt}\textbf{#1}\hfill #2\par
  \textbf{#3}\hfill #4\par\vspace{-3pt}}

\begin{document}

%=====================================================================
%  HEADER
%=====================================================================
\begin{center}
  {\LARGE\bfseries\color{accent} Swayam Singh}\\[5pt]
  \small
  \faEnvelope~\href{mailto:singhswayam008@gmail.com}{singhswayam008@gmail.com}
  \quad
  \faPhone*~+91~9116277756
  \quad
  \href{https://github.com/SwayamInSync}{\faGithub~GitHub}
  \quad
  \href{https://www.linkedin.com/in/swayam-singh-406610213/}{\faLinkedin~LinkedIn}
  \quad
  \href{https://scholar.google.com/citations?user=clLJfm8AAAAJ}{\faGraduationCap~Google Scholar}
\end{center}

%=====================================================================
%  EDUCATION
%=====================================================================
\section{Education}
\textbf{B.Tech (Computer Science and Engineering)}\hfill 2020 -- 2024 \quad\textbar\quad CGPA: 9.56\par
\textbf{University of Allahabad}\par

%=====================================================================
%  PROFESSIONAL EXPERIENCE
%=====================================================================
\section{Professional Experience}

\entry{Research Fellow}{07/2024 -- Present}{Microsoft Research}{}
\begin{tight}
  \item Engineered an \textbf{end-to-end code autoformalization pipeline} in Rust, enabling automated translation of natural language and informal code specifications into formally \textbf{verifiable representations}.
  \item Developing \textbf{foundational algorithms} that exploit \textbf{reinforcement learning gradient optimization dynamics} to improve representation efficiency, stability, and generalization in large-scale models.
  \item Developed an LLM performing on par with GPT-4o, focused on code generation, instruction-following, and reasoning; worked on post-pretraining and online/offline RL methods.
  \item Developed a new supervised fine-tuning method, \textbf{SeleKT}, that preserves generalization while significantly improving downstream task performance.
  \item \textbf{Led large-scale model training on multi-GPU clusters}, optimizing distributed training workflows and fine-tuning for high-performance code synthesis and execution.
\end{tight}

\entry{Maintainer at NumPy}{07/2024 -- Present}{NumPy}{}
\begin{tight}
  \item Review, merge, and evaluate core PRs affecting dtype architecture, numerical semantics, and cross-platform behavior.
  \item Maintain core NumPy with a focus on dtype infrastructure, numerical correctness, and long-term maintenance.
  \item Work closely with contributors and maintainers across architectures to address portability, performance, and ABI concerns.
\end{tight}

\entry{Open Source Intern -- Numerical Types \& Precision (NumPy OSS)}{07/2024 -- 09/2024}{Quansight}{}
\begin{tight}
  \item Led development of \textbf{NumPy-QuadDType}, a cross-platform 128-bit floating-point system providing consistent quadruple-precision arithmetic beyond platform-dependent long double.
  \item Implemented full NumPy dtype ecosystem integration: \textbf{casting, ufunc dispatch, scalar types, memory model, and serialization}.
  \item Built robust multi-platform distribution and testing infrastructure, handling \textbf{compiler toolchains, SIMD availability, FMA support, and architecture-specific numerical behavior}.
\end{tight}

\entry{Machine Learning Engineer Intern}{05/2022 -- 10/2022}{dataX.ai}{}
\begin{tight}
  \item Developed deep-learning models for business-market growth across Computer Vision and language modelling.
  \item Engineered a comprehensive API to automate conversion of a pretrained model into ONNX format and its seamless deployment via NVIDIA Triton Server, culminating in a 12\% reduction in VM load.
  \item Developed custom CUDA kernels to accelerate 3D image processing for medical scans, achieving a 2x speedup in segmentation tasks over existing GPU-based solutions and improving overall throughput.
\end{tight}

%=====================================================================
%  RESEARCH AND PUBLICATION
%=====================================================================
\section{Research and Publication}
\begin{tight}
  \item \textbf{Swayam Singh} et al. ``\textbf{NextCoder: Robust Learning of Diverse Code Edits}'' arXiv:2503.03656, 2025.
  \begin{tight}
    \item Published in \textbf{ICML 2025} main conference \quad\textbar\quad \textbf{ICLR 2025 DL4Code} workshop.
  \end{tight}
  \item \textbf{Swayam Singh}, Niklas Muennighoff, et al. ``\textbf{OctoPack: Instruction Tuning Code Large Language Models}'' arXiv:2308.07124, 2023.
  \begin{tight}
    \item \textbf{Spotlight (top 5\%)} at \textbf{ICLR 2024} \quad\textbar\quad \textbf{Instruction Workshop} at \textbf{NeurIPS 2023}.
  \end{tight}
  \item \textbf{Swayam Singh}, Harm de Vries, et al. ``\textbf{StarCoder: may the source be with you!}'' arXiv:2305.06161, 2023.
  \begin{tight}
    \item Published at \textbf{TMLR} (Transactions on Machine Learning Research).
  \end{tight}
\end{tight}

%=====================================================================
%  PROJECTS
%=====================================================================
\section{Projects}

\textbf{Virtual Clothing Assistant}~~{\small[PyTorch, Flask, OpenCV]}\hfill 500+ {\color{Dandelion}\faStar}~on GitHub\par
\begin{tight}
  \item Lets users try on different clothes virtually, without going to a trial room.
  \item Used \textbf{ResNet101} and \textbf{UNet} to segment the cloth and model respectively, with pose estimation via \textbf{OpenPose}; final predictions use the PyTorch implementation of \textbf{VITON}.
  \item Achieved an \textbf{SSIM} (Structural Similarity Index Metric) of 0.895.
\end{tight}

\vspace{1.5pt}
\textbf{Numpy-QuadDType}~~{\small A cross-platform Quad (128-bit) float data-type for NumPy.}\par
\begin{tight}
  \item Core developer and maintainer of a cross-platform 128-bit floating-point data type for NumPy.
  \item Designed a portable alternative to long double, eliminating architecture- and compiler-dependent irregularities.
  \item Integrated with NumPy's dtype, casting, and dispatch mechanisms while maintaining performance and correctness.
\end{tight}

\vspace{1.5pt}
\textbf{MIRA -- Multimodal Image Reconstruction with Attention}\par
\begin{tight}
  \item Designed a transformer-based multimodal architecture for single-view text/image-to-3D reconstruction using ViT encoders and triplane decoders with cross-attention.
  \item Enabled fast 3D mesh and video generation from a single image in \textless10 seconds on an A100 GPU, significantly outperforming traditional NeRF pipelines.
  \item Integrated diffusion-based image generation (SDXL) with structured 3D reconstruction for end-to-end multimodal scene synthesis.
\end{tight}

%=====================================================================
%  SKILLS
%=====================================================================
\section{Skills}
\begin{tight}
  \item Systems programming, numerical computing, cross-platform portability.
  \item C/C++, Python (NumPy internals), memory management, performance optimization.
  \item Large-scale open-source development, code review, and long-term maintenance.
\end{tight}

\end{document}
